% \iffalse meta-comment
%
% Copyright (c) 2021 Daniel Schmitz
% 
% Permission is hereby granted, free of charge, to any person obtaining a copy of this software and
% associated documentation files (the "Software"), to deal in the Software without restriction,
% including without limitation the rights to use, copy, modify, merge, publish, distribute,
% sublicense, and/or sell copies of the Software, and to permit persons to whom the Software is
% furnished to do so, subject to the following conditions:
% 
% The above copyright notice and this permission notice shall be included in all copies or substantial
% portions of the Software.
% 
% THE SOFTWARE IS PROVIDED "AS IS", WITHOUT WARRANTY OF ANY KIND, EXPRESS OR IMPLIED, INCLUDING BUT
% NOT LIMITED TO THE WARRANTIES OF MERCHANTABILITY, FITNESS FOR A PARTICULAR PURPOSE AND
% NONINFRINGEMENT. IN NO EVENT SHALL THE AUTHORS OR COPYRIGHT HOLDERS BE LIABLE FOR ANY CLAIM,
% DAMAGES OR OTHER LIABILITY, WHETHER IN AN ACTION OF CONTRACT, TORT OR OTHERWISE, ARISING FROM, OUT
% OF OR IN CONNECTION WITH THE SOFTWARE OR THE USE OR OTHER DEALINGS IN THE SOFTWARE.
%
% \fi

% \iffalse

%<*driver>
\documentclass[full]{l3doc}
\usepackage[height=1.5\baselineskip]{mahjong}
\usepackage{tabularx}
\usepackage{array}
\usepackage{booktabs}
\usepackage{multicol}
\usepackage{cleveref}
\renewcommand{\arraystretch}{1.8}
\parskip 6pt
\parindent 0pt
\EnableCrossrefs
\CodelineIndex
\RecordChanges
\begin{document}
\DocInput{mahjong.dtx} \end{document}
%</driver>
% \fi

% \CharacterTable
%  {Upper-case    \A\B\C\D\E\F\G\H\I\J\K\L\M\N\O\P\Q\R\S\T\U\V\W\X\Y\Z
%   Lower-case    \a\b\c\d\e\f\g\h\i\j\k\l\m\n\o\p\q\r\s\t\u\v\w\x\y\z
%   Digits        \0\1\2\3\4\5\6\7\8\9
%   Exclamation   \!     Double quote  \"     Hash (number) \#
%   Dollar        \$     Percent       \%     Ampersand     \&
%   Acute accent  \'     Left paren    \(     Right paren   \)
%   Asterisk      \*     Plus          \+     Comma         \,
%   Minus         \-     Point         \.     Solidus       \/
%   Colon         \:     Semicolon     \;     Less than     \<
%   Equals        \=     Greater than  \>     Question mark \?
%   Commercial at \@     Left bracket  \[     Backslash     \\
%   Right bracket \]     Circumflex    \^     Underscore    \_
%   Grave accent  \`     Left brace    \{     Vertical bar  \|
%   Right brace   \} Tilde         \~}

% \changes{v0.5}{2020/04/07}{First working version, minimal error handling}
% \changes{v0.9}{2020/04/11}{Fully functional}
% \changes{v1.0}{2020/04/14}{First complete release}
% 
% \GetFileInfo{mahjong.sty}
% \DoNotIndex{\#,\$,\%,\&,\@,\\,\{,\},\^,\_,\~,\ }
% \DoNotIndex{\advance,\begingroup,\catcode,\closein}
% \DoNotIndex{\closeout,\day,\def,\edef,\else,\empty,\endgroup}

% \title{The \textsf{mahjong} package\thanks{This document corresponds to \textsf{mahjong}~\fileversion, dated~\filedate}}
%  \author{Daniel Schmitz \\ \texttt{schmytzi@gmail.com}}
%
% \maketitle
% \tableofcontents
% \begin{abstract}
% The \textsf{mahjong} package provides a \LaTeXe{} and \LaTeX~3 interface for typesetting mahjong tiles using an extended version of MPSZ algebraic notation.
% Features include spaces, rotated, blank, and concealed tiles, as well as red fives.
% The size of the mahjong tiles can be controlled using a package option and an optional argument of |\mahjong|.
% It is primarily aimed at Riichi (aka. Japanese) Mahjong but can be used to typeset any style of mahjong.
% However, flower tiles and jokers are currently missing.
% \end{abstract}
%
% \begin{documentation}
% \section{Introduction}
% Mahjong is a tile-based game originating from China which is popular in East and South-East Asia and has since spread throughout the world.
% The \textsf{mahjong} package provides an interface for typesetting mahjong tiles and hands using MPSZ algebraic notation.
% This documentation assumes familiarity with the game in general but none of its many styles.
% Nonetheless, some basic terms will be defined because of differing vocabulary among players.
% 
% \newcommand{\rmahjong}[1]{\raisebox{-.4\height}{\mahjong{#1}}}
% \section{Mahjong Tiles}
% \subsection{Suited Tiles}
% The suited tiles are referred to as follows:
% 
% \begin{tabular}{@{} l l @{}}
% \toprule
% \textbf{Suit} & \textbf{Tiles} \\
% \midrule
% Bamboo & \rmahjong{1234506789s} \\
% Dots & \rmahjong{1234506789p} \\
% Character & \rmahjong{1234506789m} \\
% \bottomrule
% \end{tabular}
% 
% Suited tiles are referred to using the pattern \meta{value} \meta{suit}.
% For instance, \mahjong{4s} is called \textit{4~Bamboo}.
% 
% \subsection{Honor Tiles}
% This documentation refers to the seven honor tiles as follows:
%
% \begin{tabular}{@{}c c c c@{}}
% \toprule
% \multicolumn{4}{c}{\bfseries Winds} \\
% \mahjong{1z} & \mahjong{2z} & \mahjong{3z} & \mahjong{4z} \\
% East Wind (E) & South Wind (S) & West Wind (W) & North Wind (N) \\
% \midrule
% \multicolumn{4}{c}{\bfseries Dragons} \\
% \mahjong{5z} & \mahjong{6z} & \mahjong{7z} & \\
% White Dragon & Green Dragon & Red Dragon & \\
% \bottomrule
% \end{tabular}
%
% \section{MPSZ Algebraic Notation}
% \subsection{Standard Notation}
% \begin{table}
% \centering
% \caption[MPSZ reference]{MPSZ notation reference. Each tile is identified by its column's number and its row's letter.}
% \label{tab:mpszref}
% \begin{tabular}{ccccccccccc}
%    & \textbf{0} & \textbf{1} & \textbf{2} & \textbf{3} & \textbf{4} & \textbf{5} & \textbf{6} & \textbf{7} & \textbf{8} & \textbf{9} \\
%    \textbf{s} & \raisebox{-.75\baselineskip}{\rmahjong{0s}} & \rmahjong{1s} & \rmahjong{2s} & \rmahjong{3s} & \rmahjong{4s} & \rmahjong{5s} & \rmahjong{6s} & \rmahjong{7s} & \rmahjong{8s} & \rmahjong{9s} \\
%    \textbf{p} & \rmahjong{0p} & \rmahjong{1p} & \rmahjong{2p} & \rmahjong{3p} & \rmahjong{4p} & \rmahjong{5p} & \rmahjong{6p} & \rmahjong{7p} & \rmahjong{8p} & \rmahjong{9p} \\
%    \textbf{m} & \rmahjong{0m} & \rmahjong{1m} & \rmahjong{2m} & \rmahjong{3m} & \rmahjong{4m} & \rmahjong{5m} & \rmahjong{6m} & \rmahjong{7m} & \rmahjong{8m} & \rmahjong{9m} \\
%    \textbf{z} & & \rmahjong{1z} & \rmahjong{2z} & \rmahjong{3z} & \rmahjong{4z} & \rmahjong{5z} & \rmahjong{6z} & \rmahjong{7z} & &
% \end{tabular}
% \end{table}
% MPSZ notation refers to assigns each tile an identifier consisting of a digit and a letter (\cref{tab:mpszref}).
% For suited tiles, the digit corresponds to the tile's value and the letter to its suit, Bamboo (|s|), Dots (|p|) or Character (|m|).
% For instance, |2m| identifies \mahjong{2m} (2~Character).
% The only exception ot this rule are red fives, whose numeric value is 0.
% Red 5 Bamboo, for example, has identifier |0s|.
% Honor tiles are assigned the "suit" |z|, with |1z| -- |4z| corresponding to E, S, W and N, and |5z| -- |7z| to the white, green and red dragon, respectively.
% 
% Collections of tiles, such as melds or hands, are represented by concatenating the identifiers of the tiles they comprise.
% For instance, |3s4s5s| corresponds to \mahjong{3s4s5s}.
% Groups of tiles sharing the same suit can be abbreviated by omitting all but the last suit identifier.
% The above example can also be expressed as |345s|.
% Spaces are ignored and the notation is case-insensitive.
%
% \subsection{Extensions}
% \paragraph{Spaces.} Spaces can be inserted using |-|:
% |444s-567s| produces \mahjong{444s-567s}.
%
% \paragraph{Concealed Tiles.} \mahjong{X} Concealed (or face-down) tiles are represented by |X|.
% They don't need a suit identifier and don't act as one.
% |123s X 456s| and |123 X 456s| are therefore equivalent.
%
% \paragraph{Blank Tiles.} \mahjong{?} Blank or unknown tiles are represented by |?|.
% Just like concealed tiles, they don't change the current suit.
% |123s ? 456s| and |123 ? 456s| are equivalent, for instance.
%
% \paragraph{Rotation.} Inserting an apostrophe (|'|) rotates the \textit{preceeding} tile counter-clockwise.
% For instance, |6'66m| is rendered as \mahjong{6'66m}. 
% This can only be done once per tile, i.e. it is not possible to rotate them 180° or 270°.
% When you want to rotate the last tile of a group, it doesn't matter whether the apostrophe appears before or after the suit, so |77'm| and |77m'| are equivalent.
%
%  \DeleteShortVerb{\"}
% \paragraph{Rotation and Stacking.} Quotes (|"|) cause the \textit{preceeding} tile to be rendered as two rotated and stacked tiles.
% For instance, |77"7z| produces \mahjong{77"7z}.
%
% \section{Typesetting Mahjong Tiles in Your Document}
% \DescribeMacro{mahjong}
% The main interface is |\mahjong| \oarg{height} \marg{hand}.
% \meta{hand} refers to a tile sequence in MPSZ notation as discussed above and \meta{height} specifies the height of the rendered mahjong tiles.
% If \meta{height} is not specified, the default height is used.
%
% \DescribeMacro{\mahjong_typeset_hand:n} \DescribeMacro{\mahjong_typeset_hand:x}.
% The \LaTeX~3 interface for rendeering mahjong tiles are |\mahjong_typeset_hand:n| and its variants.
% This macro accepts the hand to be rendered in MPSZ notation.
% The height can be specified by setting |\l_mahjong_tile_height| \DescribeMacro{\l_mahjong_tile_height} and the default height is saved in |\g_mahjong_default_height|. \DescribeMacro{\g_mahjong_default_height}
% 
% The default height can be set using the package's |height| parameter.
% For instance, |\usepackage[2\baselineskip]{mahjong}| sets the default size of mahjong tiles to double the value of |\baselineskip| in the context they are rendered in.
% The parameter defaults to |\baselineskip|.
%
% \section{Acknowledgments}
% The mahjong tiles used in this package were created by GitHub user \href{https://github.com/FluffyStuff}{FluffyStuff}.
% The original repository is \href{https://github.com/FluffyStuff/riichi-mahjong-tiles}{FluffyStuff/riichi-mahjong-tiles}, used under CC-BY Version 4.0.
% \end{documentation}
% \begin{implementation}
%
%    \begin{macrocode}
%<*pkg>
%<@@=mahjong>
\NeedsTeXFormat{LaTeX2e}[2019/10/01]
\RequirePackage{expl3}
\ProvidesExplPackage{mahjong}{2021/04/14}{1.0}{Typeset Mahjong Hands}
\RequirePackage{xparse}
\RequirePackage{l3keys2e}
\RequirePackage{graphicx}
\RequirePackage{stackengine}
%    \end{macrocode}

%    \begin{macrocode}
\msg_new:nnnn {mahjong} {invalid_token}
{Token ~ #1 ~ is ~ not ~ valid ~ in ~ MPSZ ~ notation}
{Valid ~ tokens ~ are ~ digits ~ 0-9, ~ m, ~ p, ~ s, ~ z, ~ x, ~ -, ~ ?, ~ ', ~ and ~ " }
\msg_new:nnnn {mahjong} {unknown_tile}
{I~don't~know~tile~#1.}
{Please~check~the~documentation~for~recognized~tiles.}
\msg_new:nnnn {mahjong} {unknown_orientation}
{Orientation ~ #1 ~ is ~ unknown}
{This ~ should ~ not ~ happen.~ Please ~ create ~ a ~ bug ~ report.}

\keys_define:nn {mahjong} {
    height .dim_gset:N = \g_mahjong_default_height
}

% Identifiers for all suits
\cs_new:Npn \c_@@_suits_tl {mpsz}
% Allowed tokens
\cs_new:Npn \c_@@_allowed_tokens_tl {0123456789mpsz-?x'"}

% Variables have to be declared globally
\tl_new:N \l_@@_suit_tl
\tl_new:N \l_@@_tiles_tl
\tl_new:N \l_@@_reversed_tl
\tl_new:N \l_@@_hand_tl
\tl_new:N \l_@@_current_suit_tl
\tl_new:N \l_@@_current_group_tl
\tl_new:N \l_@@_current_char


\dim_set:Nn \g_mahjong_default_height {\baselineskip}
\dim_new:N \l_mahjong_tile_height

\int_new:N \l_@@_tile_orientation_int
\seq_new:N \l_@@_tile_maps_seq
\str_new:N \l_@@_file_path_str


\ProcessKeysPackageOptions{mahjong}
%    \end{macrocode}

% \begin{macro}[aux]{\@@_make_tile:nn, \@@_make_tile:VV, \@@_make_tile:xV, \@@_make_tile:nV}
% Inserts a mahjong tile into the input stream.
% This functions only handles that use the front background and have a forground,
% i.e. regular and blank tiles.
%    \begin{macrocode}
\cs_set:Npn \@@_make_tile:nn #1#2 {
    \file_if_exist:nTF {#1} {
        \int_case:nnF {#2} {
            {0} {
                \stackinset{c}{0pt}{c}{0pt}{
                    \includegraphics[
                        angle=0,
                        height=.85\l_mahjong_tile_height]
                    {#1}
                }{
                    \includegraphics[
                        angle=0, 
                        height=\l_mahjong_tile_height]
                    {tiles/Front.pdf}
                }
            } {1} {
                \stackinset{c}{0pt}{c}{0pt}{
                    \includegraphics[
                        angle=90,
                        width=.85\l_mahjong_tile_height]
                    {#1}
                }{
                    \includegraphics[
                        angle=90,
                        width=\l_mahjong_tile_height]
                    {tiles/Front.pdf}
                }
            } {2} {
                % Stack 2 rotated tiles on top of each other.
                \stackon [0pt] {
                    \@@_make_tile:nn {#1} {1}
                } {
                    \@@_make_tile:nn {#1} {1}
                }
            }
        } {
            \msg_fatal:nnx {mahjong} {unknown_orientation} {#2}
        }
    } {
        \msg_error:nnx {mahjong} {unknown_tile} {#1}
    }
}

\cs_generate_variant:Nn \@@_make_tile:nn {VV, xV, nV}
%    \end{macrocode}
% \end{macro}
% \begin{macro}{\mahjong_typeset_hand:n, \mahjong_typeset_hand:x}
% Parseses and typesets a mahjong hand in MPSZ notation.
% Set |\l_mahjong_tile_height| to control the tiles' size.
%    \begin{macrocode}
% Parses a full hand
\cs_set:Npn \mahjong_typeset_hand:n #1 {
    \tl_set:Nx \l_@@_hand_tl {\text_lowercase:n {#1}} 
    % MPSZ notation is easier to parse right-to-left, so reverse string
    % There is no string reversal function but we can reverse token lists
    % Token lists and strings are freely convertible between each other
    \tl_set:Nx \l_@@_reversed_tl {\tl_reverse:V \l_@@_hand_tl} 
    \tl_map_variable:NNn \l_@@_reversed_tl \l_@@_current_char {
        % Check if we recognize the current token
        \exp_args:NVV \tl_if_in:nnF \c_@@_allowed_tokens_tl \l_@@_current_char {
            \msg_error:nnx {mahjong} {invalid_token} {\l_@@_current_char}
        }
        \exp_args:NVV \tl_if_in:nnTF \c_@@_suits_tl \l_@@_current_char {
            % If we've found a suit identifier, change the current suit
            \tl_set:NV \l_@@_current_suit_tl \l_@@_current_char
        } {
            \str_case:VnF \l_@@_current_char {
                {'} {
                    % An apostrophe indicates that the next tile is rotated
                    \int_set:Nn \l_@@_tile_orientation_int {1}
                }
                {"} {
                    % Quotes mean the next tile is actually 2 rotated tiles stacked on top of each other
                    \int_set:Nn \l_@@_tile_orientation_int {2}
                }
            } {
                % Default case: We've got a complete tile identifier
                % Turn it into a property list
                \prop_clear:N \l_tmpa_prop
                \prop_put:NnV \l_tmpa_prop {suit} \l_@@_current_suit_tl
                \prop_put:NnV \l_tmpa_prop {face} \l_@@_current_char
                \prop_put:NnV \l_tmpa_prop {orientation} \l_@@_tile_orientation_int
                % Add it to the beginning of the sequence (we are parsing in reverse)
                \seq_put_left:NV \l_@@_tile_maps_seq \l_tmpa_prop
                \int_set:Nn \l_@@_tile_orientation_int {0}
            }
        }
    }
    % Typesetting begins here. Sequence is in correct order
    \seq_map_variable:NNn \l_@@_tile_maps_seq \l_tmpa_prop {
        \prop_get:NnN \l_tmpa_prop {face} \l_tmpa_tl
        \prop_get:NnN \l_tmpa_prop {orientation} \l_tmpa_int
        \str_case:VnF \l_tmpa_tl {
            {-} {
                % If the current face is a dash, insert a space
                \includegraphics[height=\l_mahjong_tile_height]{tiles/space.pdf}
            } {x} {
                % Insert a flipped tile
                \int_case:nn {\l_tmpa_int} {
                    {0} { % Upright
                        \includegraphics[
                        angle=0,
                        height=\l_mahjong_tile_height]
                        {tiles/Back.pdf}
                    } {1} { % Rotated
                        \includegraphics[
                        angle=90,
                        width=\l_mahjong_tile_height]
                        {tiles/Back.pdf}
                    } {2} { % Rotated and stacked
                        \stackon [0pt] {
                            \includegraphics[
                                angle=90,
                                width=\l_mahjong_tile_height]
                                {tiles/Back.pdf}
                        } {
                            \includegraphics[angle=90,
                            width=\l_mahjong_tile_height]
                            {tiles/Back.pdf}
                        }
                    }
                }
            } {?} { % Blank tile
                \@@_make_tile:nV {tiles/Blank.pdf} \l_tmpa_int
            }
        } { % Any other tile identified by a code.
            \@@_make_tile:xV {tiles/\l_tmpa_tl\prop_item:Nn \l_tmpa_prop {suit}.pdf} \l_tmpa_int
        }
    }
    % Clear the list for the next invocation
    \seq_clear:N \l_@@_tile_maps_seq
}

% Have TeX automatically expand the argument for us
\cs_generate_variant:Nn \@@_typeset_hand:n {x}
%    \end{macrocode}
% \end{macro}
% 
% \begin{macro}{\mahjong}
% This is the only \LaTeXe{} in this package.
% It typesets a mahjong hand.
%    \begin{macrocode}
\NewDocumentCommand{\mahjong}{O{\g_mahjong_default_height} m}{
    \dim_set:Nn \l_mahjong_tile_height {#1}
    \mahjong_typeset_hand:n {#2}
}
%</pkg>
%    \end{macrocode}
% \end{macro}
% 
% \end{implementation}
%
% \PrintChanges
\endinput
